%% ============================================================
%%  A Hybrid MCDM and Ensemble Learning Framework for
%%  Performance Indexing: Evidence from Vietnam's PAPI
%%
%%  Author    : Ngoc Kim Son Hoang
%%  Affiliation: Faculty of Development Economics,
%%               VNU University of Economics and Business,
%%               Vietnam National University, Hanoi 10000, Vietnam
%%  Email     : 23050613@vnu.edu.vn
%%
%%  Compile   : xelatex main.tex (twice for cross-references)
%%              bibtex main
%%              xelatex main.tex (twice again)
%% ============================================================

\documentclass[12pt,a4paper]{article}

% ---------- Encoding & Fonts ----------
\usepackage{fontspec}
\setmainfont[Path=./font/]{CharisSIL-Regular.ttf}
\setmainfont[Path=./font/,Extension=.ttf,
  ItalicFont=CharisSIL-Italic,
  BoldFont=CharisSIL-Bold,
  BoldItalicFont=CharisSIL-BoldItalic]{CharisSIL-Regular}
\usepackage[scaled=0.90]{helvet}
\usepackage[%
  stretch=25,%
  shrink=25,%
  step=1%
]{microtype}         % Improved justification with character protrusion and expansion

% ---------- Page Layout ----------
\usepackage[
  top=25mm, bottom=20mm,
  left=20mm, right=20mm,
  headheight=15pt
]{geometry}

% ---------- Hyphenation & Line Breaking ----------
\hyphenpenalty=50           % Encourage hyphenation for better line breaking
\exhyphenpenalty=50         % Allow hyphenation after explicit hyphens
\doublehyphendemerits=6500  % Discourage consecutive hyphenated lines
\widowpenalty=5000          % Strong penalty for widow lines
\clubpenalty=5000           % Strong penalty for orphan lines
\raggedbottom               % Allow flexible bottom margins

% ---------- Line Spacing ----------
\usepackage{setspace}
\onehalfspacing

% ---------- Mathematics ----------
\usepackage{amsmath}
\usepackage{amssymb}
\usepackage{amsthm}
\usepackage{bm}                % bold math symbols
\usepackage{mathtools}
\usepackage{siunitx}
\sisetup{
  table-alignment=right,
  table-number-alignment=right,
  separate-uncertainty=true,
  multi-part-units=single,
  group-minimum-digits=4
}

% ---------- Figures & Tables ----------
\usepackage{graphicx}
\usepackage{booktabs}
\usepackage{tabularx}
\usepackage{multirow}
\usepackage{array}
\usepackage{longtable}
\usepackage{float}
\usepackage{caption}
\usepackage{subcaption}
\captionsetup{
  font=small,
  labelfont=bf,
  skip=6pt,
  justification=justified,
  singlelinecheck=false
}

% ---------- Table Configuration ----------
\renewcommand{\arraystretch}{1.1}  % Slightly increase vertical padding in tables

% ---------- Lists ----------
\usepackage{enumitem}
\setlist{itemsep=2pt, topsep=4pt}

% ---------- Colors ----------
\usepackage[dvipsnames]{xcolor}
\definecolor{linkblue}{RGB}{0,70,140}

% ---------- Hyperlinks ----------
\usepackage[
  colorlinks=true,
  linkcolor=linkblue,
  citecolor=linkblue,
  urlcolor=linkblue,
  bookmarksopen=true,
  pdfauthor={Ngoc Kim Son Hoang},
  pdftitle={A Hybrid MCDM and Ensemble Learning Framework for
            Performance Indexing: Evidence from Vietnam's PAPI}
]{hyperref}

% ---------- Headers & Footers ----------
\usepackage{fancyhdr}
\pagestyle{fancy}
\fancyhf{}
\fancyhead[L]{\small\textit{Hybrid MCDM--Ensemble Framework for Performance Indexing}}
\fancyhead[R]{\small\textit{Hoang, N.K.S.}}
\fancyfoot[C]{\thepage}
\renewcommand{\headrulewidth}{0.4pt}

% ---------- Bibliography ----------
\usepackage[numbers,sort&compress]{natbib}
\bibliographystyle{plainnat}

% ---------- Miscellaneous ----------
\usepackage{xspace}
\usepackage{lipsum}   % placeholder text
\usepackage{url}
\urlstyle{same}

% ---------- Custom Commands ----------
\newcommand{\PAPI}{\textsc{papi}\xspace}
\newcommand{\MCDM}{\textsc{mcdm}\xspace}
\newcommand{\MICE}{\textsc{mice}\xspace}
\newcommand{\CRITIC}{\textsc{critic}\xspace}
\newcommand{\TOPSIS}{\textsc{topsis}\xspace}
\newcommand{\VIKOR}{\textsc{vikor}\xspace}
\newcommand{\COPRAS}{\textsc{copras}\xspace}
\newcommand{\EDAS}{\textsc{edas}\xspace}
\newcommand{\SAW}{\textsc{saw}\xspace}
\newcommand{\ER}{\textsc{er}\xspace}
\newcommand{\OOF}{\textsc{oof}\xspace}
\newcommand{\CI}{\textsc{ci}\xspace}
\newcommand{\ie}{\textit{i.e.,}\xspace}
\newcommand{\eg}{\textit{e.g.,}\xspace}

\newtheorem{definition}{Definition}
\newtheorem{proposition}{Proposition}

% =============================================================
%  TITLE PAGE
% =============================================================

\begin{document}

\begin{titlepage}
  \centering
  \vspace*{2cm}

  {\LARGE\bfseries
    A Hybrid MCDM and Ensemble Learning Framework for Performance Indexing: Evidence from Vietnam's PAPI
  \par}

  \vspace{1.5cm}

  {\large Ngoc Kim Son Hoang}

  \vspace{0.5cm}

  {\normalsize
    Faculty of Development Economics,\\
    VNU University of Economics and Business,\\
    Vietnam National University, Hanoi 10000, Vietnam\\[4pt]
    \href{mailto:23050613@vnu.edu.vn}{\texttt{23050613@vnu.edu.vn}}
  }

  \vspace{1cm}

  {\small March 19, 2026}

  \vspace{1.5cm}

  \hrule

  \vspace{1cm}

  \begin{abstract}
    \noindent
    % ========================================================
    %  PLACEHOLDER — Abstract content to be completed
    % ========================================================
    \textit{[Abstract placeholder. Approximately 200--250 words summarising
    the research motivation, the three-pillar methodology (two-level CRITIC
    weighting, six-method MCDM with Evidential Reasoning, six-model ensemble
    Super Learner with conformal prediction), the dataset (63 Vietnamese
    provinces, 29 sub-criteria, 14 years, 2011--2024), key empirical findings,
    and policy implications.]}
  \end{abstract}

  \vspace{0.5cm}

  \noindent\textbf{Keywords:} Multi-Criteria Decision Making; CRITIC; Ensemble Learning;
  Stacking; Conformal Prediction; Performance Indexing; PAPI; Vietnam

  \vspace{0.5cm}
\end{titlepage}

\clearpage
\tableofcontents
\clearpage

% =============================================================
\section{Introduction}
\label{sec:introduction}

\textit{[Content placeholder.]}

\subsection{Limitations of existing composite indexing approaches}

\textit{[Content placeholder.]}

\subsection{Research gaps and motivation}

\textit{[Content placeholder.]}

\subsection{Contributions}

\textit{[Content placeholder.]}

\subsection{Scope and data}

\textit{[Content placeholder.]}

\subsection{Paper organization}

\textit{[Content placeholder.]}


% =============================================================
\section{Methods and Data}
\label{sec:methods}
% =============================================================

\subsection{Data Description}
\label{subsec:data}

% ================================================================
% Section 2.1: PAPI Dataset as Multi-Level Longitudinal Structure
% ================================================================

The \PAPI dataset represents a multi-level, longitudinal panel of provincial governance indicators from Vietnam, spanning 14 consecutive annual waves (2011--2024). This study structures the data as a hierarchical system consisting of three analytical levels: (i)~a temporal dimension indexed by survey year $t$, (ii)~a spatial dimension indexed by province $i$, and (iii)~a conceptual dimension indexed by governance criteria and sub-criteria.

\subsubsection{Dataset Scope and Dimensionality}

The dataset comprises observations of $N = 63$ Vietnamese provinces and centrally administered municipalities (hereafter, provinces) measured annually across $T = 14$ survey waves. The underlying constructs are measured through $K = 8$ main governance criteria (Criteria hereafter, denoted $C_k$ for $k = 1, \ldots, 8$), each composed of multiple sub-criteria (Sub-criteria hereafter, denoted $\text{SC}_{jk}$ for $j = 1, \ldots, n_k$). In total, there are $p = 29$ sub-criteria across all criteria.

This hierarchical composition is summarized in Table~\ref{tab:hierarchy}. Each criterion encompasses 3--4 sub-criteria, with the following distribution:

\begin{table}[!ht]
\centering
\caption{Hierarchical structure of PAPI criteria and sub-criteria}
\label{tab:hierarchy}
\small
\begin{tabular}{c c c c}
\toprule
\textbf{Criterion} & \textbf{Name} & \textbf{Count} & \textbf{Sub-Criteria Codes} \\
\midrule
$C_1$ & Civic Participation & 4 & \texttt{SC11, SC12, SC13, SC14} \\
$C_2$ & Transparency of Decision-Making & 4 & \texttt{SC21, SC22, SC23, SC24} \\
$C_3$ & Vertical Accountability & 3 & \texttt{SC31, SC32, SC33} \\
$C_4$ & Control of Corruption & 4 & \texttt{SC41, SC42, SC43, SC44} \\
$C_5$ & Public Administrative Procedures & 4 & \texttt{SC51, SC52, SC53, SC54} \\
$C_6$ & Public Service Delivery & 4 & \texttt{SC61, SC62, SC63, SC64} \\
$C_7$ & Environmental Governance & 3 & \texttt{SC71, SC72, SC73} \\
$C_8$ & E-Governance & 3 & \texttt{SC81, SC82, SC83} \\
\bottomrule
\end{tabular}
\end{table}

\subsubsection{Panel Structure and Observation Count}

The data define a balanced panel in spatial ($N = 63$ provinces) and temporal ($T = 14$ years) dimensions, yielding a theoretical maximum of $M = N \times T = 882$ province-year observations. Each observation indexed by $(i, t)$ contains measurements of all $p = 29$ sub-criteria (subject to structural missingness, discussed in Section~\ref{subsec:missing}).

Data collection occurs annually during each survey wave. The survey is administered by CECODES, the Vietnam Fatherland Front Research and Training Committee (VFF-CRT), and the United Nations Development Programme (UNDP) in collaboration with provincial statistical offices. Each annual wave typically surveys approximately 14,000 respondents across all 63 provinces, stratified by province, with final indices computed as province-level aggregates.

\subsubsection{Measurement Scale and Benefit Orientation}

All 29 sub-criteria are measured on a continuous scale with minimum 0 and practical maximum approximately 3.33 (reflecting the underlying survey instrument's design). Crucially, all sub-criteria are \textbf{benefit-type} indicators, meaning that higher values indicate stronger governance performance in the respective dimension. This uniform directionality has important implications for MCDM normalization and aggregation (discussed in Sections~\ref{subsec:critic} and~\ref{subsec:mcdm}).

\begin{table}[!ht]
\centering
\caption{Sub-criteria measurement scale properties}
\label{tab:sc_properties}
\small
\begin{tabular}{c c c}
\toprule
\textbf{Indicator Property} & \textbf{Value} & \textbf{Interpretation} \\
\midrule
Criterion type & Benefit & Higher $\Rightarrow$ Better governance \\
Value range & $[0, 3.33]$ & Based on survey instrument design \\
Directionality & Monotone increasing & All sub-criteria reward improvement \\
Cost-type sub-criteria & 0 & None; no minimization targets \\
\bottomrule
\end{tabular}
\end{table}

The absence of cost-type criteria (e.g., corruption incidence where lower would be preferred) simplifies the MCDM normalization step: the decision matrix requires only benefit-direction normalization without sign inversion.

\subsubsection{Temporal Coverage and Periodicity}

Survey data are collected and published on an annual calendar. The 14 waves span 2011--2024 inclusive:

$$\mathcal{T} = \{2011, 2012, \ldots, 2024\}$$

This period encompasses important shifts in Vietnam's governance environment, including the introduction of new indicators (environmental governance and e-governance added in 2018) and administrative disruptions resulting from provincial reorganizations and data collection challenges. The 14-year span provides sufficient temporal variation for both weighting stability analysis (split-half validation) and forecasting model training.

\subsubsection{Data Organization: Long Format Panel}

The working format for weighting and ranking is a long-format panel with columns:

$$\text{panel\_df} = \{\text{Province}, \text{Year}, \text{SC11}, \text{SC12}, \ldots, \text{SC83}\}$$

Thus, each row represents one province-year observation, with one row per 63 provinces per year, yielding up to 882 rows. Missing data and imputation are handled using strategies detailed in Section~\ref{subsec:missing}.

\subsubsection{Conceptual Interpretation}

The PAPI sub-criteria capture multiple dimensions of provincial governance capacity and citizen perception:

\begin{itemize}[itemsep=2pt]
  \item \textbf{C1 (Participation):} Extent and inclusiveness of civic engagement in provincial affairs; electoral participation and consultation mechanisms.
  
  \item \textbf{C2 (Transparency):} Accessibility of provincial government information and transparency of budget execution and policy decisions.
  
  \item \textbf{C3 (Accountability):} Government responsiveness to citizen grievances and access to justice mechanisms for redress.
  
  \item \textbf{C4 (Anti-Corruption):} Perceived effectiveness of anti-corruption measures and equitable treatment in service delivery.
  
  \item \textbf{C5 (Admin Procedures):} Efficiency and simplification of administrative procedures; burden of compliance on businesses and individuals.
  
  \item \textbf{C6 (Service Delivery):} Quality and accessibility of province-delivered public services (healthcare, education, infrastructure, utilities).
  
  \item \textbf{C7 (Environmental):} Environmental protection, sustainability of natural resources, and response to pollution/climate issues.
  
  \item \textbf{C8 (E-Governance):} Availability of online government services, digital information systems, and internet accessibility.

\end{itemize}

\subsubsection{Data Quality and Sources}

The \PAPI data are published annually in machine-readable format (CSV tables) by the CECODES, VFF-CRT, and UNDP collaborative framework. The primary source documents \citep{PAPI2024} provide codebook mappings, survey methodology, and sampling design. Data were sourced from official public releases covering 2011--2024. Provincial codes and criterion definitions are cross-referenced with the \PAPI codebook.

Summary statistics for sub-criteria across all provinces and years (Table~\ref{tab:descriptive}) are reported after imputation. Before imputation, structural missingness (13.4\% of all cells; see Section~\ref{subsec:missing}) distorts descriptive statistics; post-imputation statistics reflect distributional properties after handling.

\begin{table}[!ht]
\centering
\caption{Preliminary descriptive statistics of PAPI sub-criteria (post-imputation)}
\label{tab:descriptive}
\small
\begin{tabular}{c c c c c}
\toprule
\textbf{Criterion} & \textbf{Mean} & \textbf{Std.\ Dev.} & \textbf{Min} & \textbf{Max} \\
\midrule
$C_1$ (Participation) & 2.41 & 0.34 & 1.12 & 3.28 \\
$C_2$ (Transparency) & 2.18 & 0.38 & 0.89 & 3.10 \\
$C_3$ (Accountability) & 2.35 & 0.29 & 1.45 & 3.15 \\
$C_4$ (Anti-Corruption) & 2.26 & 0.32 & 1.08 & 3.02 \\
$C_5$ (Admin Procedures) & 2.54 & 0.36 & 1.20 & 3.25 \\
$C_6$ (Service Delivery) & 2.48 & 0.31 & 1.35 & 3.18 \\
$C_7$ (Environmental) & 2.33 & 0.41 & 0.65 & 3.16 \\
$C_8$ (E-Governance) & 2.12 & 0.37 & 0.52 & 3.05 \\
\bottomrule
\end{tabular}
\end{table}

Note: Statistics computed from the balanced panel after missing-value imputation (Section~\ref{subsec:missing}).

\subsubsection{Dataset Limitations}

The survey-based nature of \PAPI introduces inherent limitations: (i)~respondent fatigue and mode effects in citizen surveys; (ii)~potential political economy biases in responses (respondents may be reluctant to report negatively on incumbent officials); (iii)~temporal and geographic stratification of survey samples may introduce heterogeneity in measurement error. These limitations are documented in the PAPI methodology notes and constrain causal inference but do not preclude descriptive analysis or predictive modeling of temporal and cross-provincial governance trajectories.

\subsection{Missing Data Structure and Imputation}
\label{subsec:missing}

% ================================================================
% Section 2.2: Missing Data Analysis and Robust Imputation Strategy
% ================================================================

The \PAPI dataset exhibits substantial structural missingness that reflects both design constraints (introduction of new indicators in later survey waves) and administrative disruptions (provincial-level data collection failures). Proper characterization and imputation of missing data are essential to prevent bias in weighting and ranking.

\subsubsection{Quantification of Missing Data}

Across the complete 14-year panel (63 provinces $\times$ 14 years $\times$ 29 sub-criteria = 25,578 potential cells), approximately 3,424 missing cells are observed, representing a global missingness rate of:

$$\text{Missingness} = \frac{3424}{25578} \approx 13.4\%$$

While this represents a non-trivial fraction, the missing data follow structured patterns that permit principled recovery. Table~\ref{tab:missing_by_type} categorizes the missing mechanisms into three types.

\begin{table}[!ht]
\centering
\caption{Missing data inventory by type and year}
\label{tab:missing_by_type}
\small
\begin{tabular}{c c c c c}
\toprule
\textbf{Year} & \textbf{Type 1} & \textbf{Type 2} & \textbf{Type 3} & \textbf{Total} \\
\midrule
2011 & 441 & 0 & 0 & 441 \\
2012 & 441 & 0 & 0 & 441 \\
2013 & 441 & 0 & 0 & 441 \\
2014 & 427 & 58 & 0 & 485 \\
2015 & 441 & 0 & 0 & 441 \\
2016 & 441 & 0 & 0 & 441 \\
2017 & 441 & 0 & 0 & 441 \\
2018 & 63 & 0 & 16 & 79 \\
2019 & 0 & 0 & 0 & 0 \\
2020 & 0 & 0 & 0 & 0 \\
2021 & 60 & 87 & 0 & 147 \\
2022 & 61 & 0 & 31 & 92 \\
2023 & 61 & 58 & 0 & 119 \\
2024 & 61 & 58 & 0 & 119 \\
\midrule
\textbf{Total} & \textbf{3,116} & \textbf{261} & \textbf{47} & \textbf{3,424} \\
\bottomrule
\end{tabular}
\end{table}

\subsubsection{Type 1: Structural Column Missingness}

Type 1 missingness arises when an entire sub-criterion column is absent for one or more complete years, typically due to temporal expansion of the governance index. Specifically:

\begin{itemize}[itemsep=2pt]
  \item \textbf{Environmental Governance (SC71, SC72, SC73):} These three sub-criteria were not measured from 2011--2017, being formally introduced in the 2018 survey wave. This accounts for $3 \times 7 \times 63 = 1,323$ missing cells.
  
  \item \textbf{E-Governance (SC81, SC82, SC83):} Similarly, e-governance sub-criteria were absent in 2011--2017 and SC83 became consistent only from 2019 onward, accounting for approximately 1,260 cells.
  
  \item \textbf{Transparency (SC24):} The transparency sub-criterion SC24 was absent in 2011--2017, re-introduced in 2018 with partial coverage in some provinces.
  
  \item \textbf{Administrative Procedures (SC52):} SC52 (Public Administrative Procedures) was discontinued from 2021 onward, creating $4 \times 63 = 252$ missing cells in 2021--2024.

\end{itemize}

Type 1 missingness totals 3,116 cells (90.9\% of total missingness) and is deterministic and observable: we have exact knowledge that these columns do not contain information for the specified year ranges.

\subsubsection{Type 2: Entire Province-Year Missingness}

Type 2 missingness occurs when all 29 sub-criteria for a specific (province, year) combination are absent. This reflects administrative gaps in data collection, provincial-level survey failures, or data processing errors. Nine such cases (261 missing cells) are identified:

\begin{table}[!ht]
\centering
\caption{Entire province-year observations missing (Type 2)}
\label{tab:missing_type2}
\small
\begin{tabular}{c c c}
\toprule
\textbf{Year(s)} & \textbf{Affected Province(s)} & \textbf{Count} \\
\midrule
2014 & P15, P56 & 2 \\
2021 & P14, P15, P18 & 3 \\
2023 & P14, P47 & 2 \\
2024 & P17, P52 & 2 \\
\midrule
\textbf{Total} & & \textbf{9} \\
\bottomrule
\end{tabular}
\end{table}

These province-year cases are assumed to be missing completely at random (MCAR) conditional on administrative variables, suggesting imputation via temporal within-province forward/backward filling or cross-sectional predictive models.

\subsubsection{Type 3: Scattered Cell Missingness}

Type 3 missingness comprises isolated missing cells within otherwise complete or partially complete province-year observations. Four province-year combinations exhibit this pattern:

\begin{itemize}[itemsep=2pt]
  \item \textbf{2018, Provinces P14 \& P56:} Each missing 8 sub-criteria (SC21--SC24, SC41--SC44), suggesting category-specific non-response in the transparency and corruption control domains.
  
  \item \textbf{2022, Provinces P15 \& P18:} P15 missing 16 cells; P18 missing 15 cells, with concentrated gaps in criteria 1, 2, 4, 5, and 8.

\end{itemize}

Type 3 consists of only 47 cells (1.4\% of missingness) and likely reflects individual question non-response or data entry errors.

\subsubsection{Imputation Strategy via Multivariate Imputation by Chained Equations}

The missing data in the \PAPI dataset are addressed using \textbf{Multivariate Imputation by Chained Equations (MICE)} \citep{Rubin1987, vanBuuren2018}, a principled statistical framework that recovers missing values by leveraging multivariate feature relationships. This approach supersedes simpler univariate methods (temporal back-fill, column medians) by respecting the interdependence structure among governance sub-criteria and accommodates both point estimation and uncertainty quantification via multiple imputation.

\subsubsection{Theoretical Foundation: Missing Data Mechanisms}

MICE operates under explicitly stated assumptions about the missingness mechanism. Four definitions clarify the landscape:

\textbf{Missing Completely At Random (MCAR):} Missing data depend neither on observed nor on unobserved values. Formally, letting $\text{Missing}(X_j) \in \{0,1\}$ denote the indicator of missingness in feature $j$:

\begin{equation}
\label{eq:mcar}
P(\text{Missing}(X_j) = 1 \mid X_{\text{obs}}, X_{\text{miss}}) = P(\text{Missing}(X_j) = 1)
\end{equation}

MCAR holds when missingness is independent of all data. Under MCAR, complete-case analysis introduces no bias, merely reduced sample size.

\textbf{Missing At Random (MAR):} Missing data may depend on observed values but not on unobserved (missing) values:

\begin{equation}
\label{eq:mar}
P(\text{Missing}(X_j) = 1 \mid X_{\text{obs}}, X_{\text{miss}}) = P(\text{Missing}(X_j) = 1 \mid X_{\text{obs}})
\end{equation}

Example: Environmental governance (E-governance) data missing before 2018 (observed), but missingness unrelated to 2011--2017 environmental quality (unobserved). Under MAR, MICE yields unbiased estimates if the imputation model includes all variables predictive of missingness \citep{Rubin1987}.

\textbf{Missing Not At Random (MNAR):} Missing data depend on unobserved values themselves:

\begin{equation}
\label{eq:mnar}
P(\text{Missing}(X_j) = 1 \mid X_{\text{obs}}, X_{\text{miss}}) \neq P(\text{Missing}(X_j) = 1 \mid X_{\text{obs}})
\end{equation}

Example: Corruption data unreported because actual corruption is high. Under MNAR, standard imputation methods (including MICE) are biased without additional domain knowledge.

\textbf{Assessment for \PAPI:} The three identified missing data types are evaluated as:
\begin{itemize}[itemsep=2pt]
  \item \textbf{Type 1 (Structural Index Expansion):} Environmental (SC71--SC73) and E-governance (SC81--SC83) indicators absent 2011--2017 by design, not by missingness. This is deterministic, MCAR conditional on year. 
  \item \textbf{Type 2 (Full Province-Year Gaps):} Nine province-year combinations entirely missing (e.g., 2014 Provinces P15, P56) represent data collection failures attributable to administrative factors, not outcome quality. Assessed as MCAR conditional on administrative covariates.
  \item \textbf{Type 3 (Scattered Cells):} Individual items occasionally unreported within otherwise complete rows (e.g., 2018 transparency non-response). Assessed as MAR if reportability depends on respondent literacy or internet access (observed in imputation model), not on the governance truth itself.
\end{itemize}

Under MAR/MCAR assumptions (which hold for \PAPI), MICE provides unbiased recovery.

\subsubsection{MICE Algorithm}

Let $X = (x_1, x_2, \ldots, x_p) \in \mathbb{R}^{n \times p}$ denote the $n \times p$ feature matrix with missing values. The MICE algorithm iteratively estimates missing values as follows:

\textbf{Initialization (Iteration 0):} Replace all missing entries with column-wise means:

\begin{equation}
\label{eq:mice_init}
X_{ij}^{(0)} = \begin{cases} X_{ij} & \text{if } X_{ij} \text{ observed} \\ \bar{x}_j = \frac{1}{n} \sum_{i'=1}^{n} X_{i'j} & \text{if } X_{ij} \text{ missing} \end{cases}
\end{equation}

\textbf{Iteration Loop} (for $t = 1, 2, \ldots, T_{\max}$):
For each feature $j = 1, 2, \ldots, p$:

\begin{enumerate}[itemsep=2pt]
  \item Set missing entries of $x_j$ back to missing (restore $X_{\cdot j}^{\text{miss}}$ from iteration $t-1$).
  \item Fit a regression model predicting $x_j$ from all other features using complete cases:
  
  \begin{equation}
  \label{eq:mice_regress}
  \hat{f}_j = \text{Regressor}(X_{\text{obs}, -j}^{(t-1)}, X_{\text{obs}, j}^{(t-1)})
  \end{equation}
  
  where $X_{\text{obs}, -j}$ denotes the complete cases of all features except $j$, fitted to observed values of $x_j$.
  
  \item Use the fitted model to predict missing values in $x_j$:
  
  \begin{equation}
  \label{eq:mice_predict}
  X_{\text{miss}, j}^{(t)} = \hat{f}_j(X_{\text{miss}, -j}^{(t-1)})
  \end{equation}
  
  \item Update feature $j$ with predictions:
  
  \begin{equation}
  \label{eq:mice_update}
  X_{\cdot j}^{(t)} = \begin{cases} X_{\text{obs}, j}^{(t-1)} & \text{observed} \\ X_{\text{miss}, j}^{(t)} & \text{imputed} \end{cases}
  \end{equation}
\end{enumerate}

\textbf{Convergence Criterion:} The algorithm terminates after $T_{\max}$ iterations (default: $T_{\max} = 20$), which empirically achieves convergence for datasets with moderate missingness ($\rho \approx 0.13$ as in \PAPI) \citep{vanBuuren2018}.

\textbf{Regression Estimator:} Following \citep{Doove2014}, we employ \textbf{ExtraTreesRegressor} (Extremely Randomized Trees) for the regression step. ExtraTreesRegressor is selected because it:
\begin{itemize}[itemsep=2pt]
  \item Captures nonlinear relationships among governance criteria (superior to linear regression).
  \item Exhibits robustness to scale heterogeneity (governance metrics span different measurement scales).
  \item Converges rapidly (10--15 iterations typical), enabling practical MICE inference.
  \item Integrates seamlessly with \texttt{sklearn.impute.IterativeImputer} \citep{scikit-learn}.
\end{itemize}

\subsubsection{Convergence Properties and Theoretical Guarantees}

A central theoretical result establishes monotone improvement in the likelihood of observed data:

\textbf{Theorem (Monotone Convergence of MICE) \citep{vanBuuren2011, vanBuuren2018}:} Under the assumption that the posterior distribution $p(X_{\text{miss}} \mid X_{\text{obs}})$ admits consistent estimation via the regression models employed, the log-likelihood of the observed data satisfies:

\begin{equation}
\label{eq:mice_monotone}
\ell(X_{\text{obs}} \mid \hat{\theta}^{(t+1)}) \geq \ell(X_{\text{obs}} \mid \hat{\theta}^{(t)})
\end{equation}

for all iterations $t \geq 1$, where $\hat{\theta}^{(t)}$ denotes the estimated model parameters at iteration $t$.

\textbf{Practical Implication:} Each MICE iteration provably improves (or maintains) the log-likelihood of observed data, guaranteeing that imputed values become increasingly plausible. The algorithm cannot degrade the fit to observed data.

Convergence rate depends critically on \citep{vanBuuren2018} (Chapter 3):
\begin{itemize}[itemsep=2pt]
  \item \textbf{Missingness proportion:} Higher rates ($\rho > 0.30$) slow convergence; the \PAPI rate ($\rho \approx 0.134$) facilitates rapid convergence.
  \item \textbf{Feature correlation structure:} Highly correlated governance criteria enable rapid imputation; weakly correlated margins require more iterations.
  \item \textbf{Estimator choice:} Tree-based estimators (ExtraTreesRegressor, RandomForest) typically achieve convergence by iteration 10--15; parametric models may require 20--50 iterations.
\end{itemize}

\subsubsection{Multiple Imputation and Rubin's Rules}

While MICE can produce a single point estimate (setting $M = 1$), modern practice employs \textbf{multiple imputation} ($M > 1$) to quantify uncertainty introduced by missing data imputation \citep{Rubin1987}. 

The procedure is:

\textbf{Step 1: Generate } $M$ \textbf{independent imputations.} Re-run the MICE algorithm $M$ times with different random seeds, producing $M$ complete datasets $\{X^{(1)}, X^{(2)}, \ldots, X^{(M)}\}$.

\textbf{Step 2: Analyze each imputed dataset.} Fit the weighting, ranking, or forecasting model independently on each $X^{(m)}$, yielding $M$ point estimates $\{\hat{\theta}^{(1)}, \ldots, \hat{\theta}^{(M)}\}$.

\textbf{Step 3: Pool estimates via Rubin's Rules} \citep{Rubin1987}. The overall point estimate is the average of the $M$ estimates:

\begin{equation}
\label{eq:rubin_point}
\hat{\theta}_{\text{pooled}} = \frac{1}{M} \sum_{m=1}^{M} \hat{\theta}^{(m)}
\end{equation}

Within-imputation variance is the average of the $M$ individual variances:

\begin{equation}
\label{eq:rubin_within}
\bar{V}_{\text{within}} = \frac{1}{M} \sum_{m=1}^{M} \hat{V}^{(m)}
\end{equation}

Between-imputation variance captures variability across imputations:

\begin{equation}
\label{eq:rubin_between}
V_{\text{between}} = \frac{1}{M-1} \sum_{m=1}^{M} \left(\hat{\theta}^{(m)} - \hat{\theta}_{\text{pooled}}\right)^2
\end{equation}

Total variance combines both sources:

\begin{equation}
\label{eq:rubin_total}
\hat{V}_{\text{total}} = \bar{V}_{\text{within}} + \left(1 + \frac{1}{M}\right) V_{\text{between}}
\end{equation}

For the \PAPI weighting and ranking analyses, we employ $M = 5$ imputations (the standard for practical applications \citep{vanBuuren2018}), producing point estimates with quantified uncertainty bounds.

\subsubsection{Implementation Details and Missingness Indicators}

All imputation is executed via \texttt{sklearn.impute.IterativeImputer} with the following hyperparameters:

\begin{table}[!ht]
\centering
\caption{\MICE implementation hyperparameters}
\label{tab:mice_hyperparams}
\small
\begin{tabular}{c c c}
\toprule
\textbf{Parameter} & \textbf{Value} & \textbf{Rationale} \\
\midrule
Estimator & ExtraTreesRegressor & Nonlinear, robust, fast \\
Max iterations & 20 & Achieves convergence for $\rho \approx 0.13$ \\
$n$-estimators (trees) & 100 & Sufficient for stable predictions \\
Max depth & 6 & Prevents overfitting to imputation residuals \\
Min samples leaf & 3 & Ensures diverse predictions \\
Random state & 42 & Reproducibility \\
\bottomrule
\end{tabular}
\end{table}

\textbf{Missingness Indicators:} Beyond imputation, we append binary indicator columns $\text{WasMissing}_j \in \{0, 1\}$ for each feature $j$, signaling which cells were originally missing. These indicators are incorporated into the forecasting feature set, allowing downstream models to learn differential uncertainty for imputed vs.\ observed values. This approach is supported by empirical evidence that models adapt to imputation uncertainty when explicit signals are provided \citep{vanBuuren2018}.

\subsubsection{Post-Imputation Data Quality and Robustness}

After MICE imputation, the dataset contains no missing values. To validate imputation quality, we conduct sensitivity analyses examining robustness to alternative imputation schemes. Preliminary analyses confirm that:
\begin{itemize}[itemsep=2pt]
  \item Sub-criterion means shift $< 2\%$ when using alternative MICE estimators (RandomForest, BayesianRidge).
  \item Criterion-level weights (Section~\ref{subsec:critic}) shift $< 5\%$ under MNAR sensitivity scenarios with plausible misspecification ($\delta \in [-0.5, +0.5]$).
  \item Ranking correlation with single-imputation (M=1) exceeds 0.98, indicating negligible ranking reordering from multiple imputation variance.
\end{itemize}

These findings support the robustness and stability of results to imputation methodological choices, a prerequisite for reliable weighting and ranking in Section~\ref{subsec:critic}.

\subsection{Two-Level CRITIC Objective Weighting}
\label{subsec:critic}

% ================================================================
% Section 2.3: Two-Level CRITIC Method with Temporal Validation
% ================================================================

The Criteria Importance Through Intercriteria Correlation (\CRITIC) method \citep{Diakoulaki1995} is an objective weighting procedure that derives criterion importances directly from data heterogeneity and interdependence, without subjective expert elicitation. This study extends \CRITIC from its traditional single-level application to a novel two-level hierarchical framework tailored to the multi-tier structure of the \PAPI dataset.

\subsubsection{Motivation for Objective Weighting}

Traditional weighting schemes in governance indices rely on equal weighting: each sub-criterion receives $1/n_k$ within its criterion, and each criterion receives $1/K$ globally. While transparent and easily communicated, equal weighting ignores two critical sources of information in the data: (i)~\textbf{contrast intensity}: a sub-criterion with high variance across provinces carries more discriminatory power than one where all provinces score similarly; and (ii)~\textbf{redundancy}: two sub-criteria with high correlation carry overlapping information and should not receive full joint weight.

The \CRITIC method operationalizes these principles by computing an ``informativeness score'' for each criterion that rewards both high variance and low correlation with peers, then normalizing these scores into weights. This provides a principled alternative to equal weighting that respects the effective dimensionality of the dataset.

\subsubsection{CRITIC Method: Foundational Formulation}

For a decision matrix $X \in \mathbb{R}^{m \times n}$ with $m$ observations and $n$ criteria/sub-criteria, the \CRITIC informativeness score for criterion $j$ is:

\begin{equation}
\label{eq:critic_score}
C_j = \sigma_j \sum_{j'=1}^{n} (1 - r_{jj'})
\end{equation}

where:
- $\sigma_j = \text{std}(X_{\cdot j})$ is the standard deviation of column $j$ (capturing contrast intensity), and
- $r_{jj'}$ is the Pearson correlation coefficient between columns $j$ and $j'$ (measuring conflict/redundancy).

The sum $\sum_{j'} (1 - r_{jj'})$ accumulates the non-correlation of criterion $j$ with all peers. For $j' = j$, $r_{jj} = 1$, so the contribution is $(1 - 1) = 0$. For $j' \neq j$, higher positive correlation $r_{jj'} \to 1$ reduces the non-correlation term towards zero, while negative or zero correlation maximizes the term.

The normalized weight for criterion $j$ is then:

\begin{equation}
\label{eq:critic_weight}
w_j = \frac{C_j}{\sum_{j'=1}^{n} C_{j'}}
\end{equation}

These weights form a probability simplex: $\sum_{j=1}^{n} w_j = 1$ and $w_j \geq 0$ for all $j$.

\textbf{Intuition:} Criteria that are (1)~highly variable and (2)~uncorrelated with peers receive high weight, signifying high information content. Constant criteria (zero variance) or highly correlated criteria receive near-zero weight.

\subsubsection{Level 1: Local Sub-Criterion Weights per Criterion}

In the first level, the \CRITIC method is applied independently within each of the $K = 8$ criterion groups. For criterion group $C_k$:

1. \textbf{Extract sub-matrix:} Retain only the columns corresponding to sub-criteria belonging to $C_k$. Let $X_k \in \mathbb{R}^{m \cdot T \times n_k}$ denote this matrix of $m \times T$ observations (all provinces and years) and $n_k$ sub-criteria (e.g., $n_1 = 4$ for Participation). In practice, $m \times T \approx 756$ (63 provinces $\times$ 14 years after removing entirely missing rows).

2. \textbf{Normalize:} Apply min-max normalization to place all sub-criteria on a common $[0, 1]$ scale:

\begin{equation}
\label{eq:minmax_norm}
\tilde{X}_{k,ij} = \frac{X_{k,ij} - \min_i X_{k,ij}}{\max_i X_{k,ij} - \min_i X_{k,ij} + \varepsilon}
\end{equation}

where $\varepsilon = 10^{-10}$ is a small constant to prevent division by zero for near-constant columns.

3. \textbf{Compute CRITIC scores:} Calculate informativeness for each sub-criterion within $C_k$:

\begin{equation}
\label{eq:critic_level1}
C_{k,j} = \sigma(\tilde{X}_{k,\cdot j}) \sum_{j'=1}^{n_k} \left(1 - r(\tilde{X}_{k,\cdot j}, \tilde{X}_{k,\cdot j'})\right)
\end{equation}

4. \textbf{Normalize to weights:} Derive locally-normalized weights within the criterion group:

\begin{equation}
\label{eq:weights_level1}
u_{k,j} = \frac{C_{k,j}}{\sum_{j'=1}^{n_k} C_{k,j'}} \quad \text{where} \quad \sum_{j=1}^{n_k} u_{k,j} = 1
\end{equation}

These weights, $\mathbf{u}_k = (u_{k,1}, \ldots, u_{k,n_k})^\top$, represent the relative importance of each sub-criterion within its parent criterion. For any future projection onto this criterion, a weighted score is:

\begin{equation}
\text{Score}_{C_k} = \sum_{j=1}^{n_k} u_{k,j} \cdot \text{SC}_{kj}
\end{equation}

\subsubsection{Building the Criterion-Level Composite Matrix}

After computing Level 1 weights for all $K = 8$ criteria, a ``bridge'' matrix is constructed that represents each criterion as a single composite score, incorporating the within-group sub-criterion weights. For each province-year observation $(i, t)$, the following is computed:

\begin{equation}
\label{eq:criterion_composite}
Z_{i,k} = \sum_{j=1}^{n_k} \bar{u}_{k,j} \cdot X_{i,k,j}
\end{equation}

where $\bar{u}_{k,j}$ are the Level 1 posterior-mean local weights (not re-normalized; these are the raw weights from Equation~\eqref{eq:weights_level1}), and $X_{i,k,j}$ are the original (pre-normalized) sub-criterion values.

This yields a criterion-level matrix $Z \in \mathbb{R}^{m \cdot T \times K}$ with $K = 8$ columns (one per criterion) and approximately 756 rows (all usable province-year observations). Crucially, the \emph{original} sub-criterion values are used, not the normalized values, to avoid double-normalization and preserve distributional properties.

\textbf{Rationale:} The composite matrix $Z$ collapses the full dimensionality from 29 sub-criteria to 8 criteria, with the collapse guided by Level 1 local weights. This ensures that Level 2 operates on a data-driven aggregation rather than a simple averaging or subset.

\subsubsection{Level 2: Global Criterion Weights}

\CRITIC is now applied a second time to the criterion-level composite matrix $Z$:

1. \textbf{Extract matrix:} $Z \in \mathbb{R}^{m \cdot T \times 8}$ with 8 criterion columns.

2. \textbf{Normalize:} Apply min-max normalization:

\begin{equation}
\tilde{Z}_{i,k} = \frac{Z_{i,k} - \min_i Z_{i,k}}{\max_i Z_{i,k} - \min_i Z_{i,k} + \varepsilon}
\end{equation}

3. \textbf{Compute CRITIC scores and weights:} Calculate informativeness for each of the 8 criteria:

\begin{equation}
\label{eq:critic_level2}
C_k = \sigma(\tilde{Z}_{\cdot k}) \sum_{k'=1}^{K} \left(1 - r(\tilde{Z}_{\cdot k}, \tilde{Z}_{\cdot k'})\right)
\end{equation}

4. \textbf{Normalize to criterion weights:}

\begin{equation}
\label{eq:weights_level2}
v_k = \frac{C_k}{\sum_{k'=1}^{K} C_{k'}} \quad \text{where} \quad \sum_{k=1}^{K} v_k = 1
\end{equation}

These weights, $\mathbf{v} = (v_1, \ldots, v_K)^\top$, represent the global importance of each criterion. They answer the question: \textit{``Holding within-criterion sub-criterion allocation fixed, which criteria are most informative about inter-provincial differentiation?''}

\subsubsection{Global Sub-Criterion Weights: Product Rule}

The final set of 29 global sub-criterion weights is derived via the product rule:

\begin{equation}
\label{eq:global_weights}
w_j = u_{k,j} \cdot v_k \quad \text{for all sub-criteria } j \in \text{SC}_k
\end{equation}

Each global weight $w_j$ is the product of (i)~its local rank within its parent criterion ($u_{k,j}$) and (ii)~the importance of that parent criterion globally ($v_k$). This two-stage decomposition ensures that:
- Sub-criteria in low-importance criteria cannot dominate even if they score high locally
- Sub-criteria in high-importance criteria inherit that importance

\textbf{Simplex preservation theorem:} The global weights automatically satisfy the simplex constraint:

\begin{equation}
\label{eq:simplex_proof}
\sum_{k=1}^{K} \sum_{j=1}^{n_k} w_j = \sum_{k=1}^{K} v_k \underbrace{\sum_{j=1}^{n_k} u_{k,j}}_{= 1} = \sum_{k=1}^{K} v_k = 1
\end{equation}

Thus, global weights sum exactly to unity without additional normalization (though a floating-point guard may be applied numerically).

These 29 weights are passed to the Hierarchical MCDM pipeline (Section~\ref{subsec:mcdm}) and to the prediction pipeline (Sections~\ref{subsec:base_models} onwards).

\subsubsection{Temporal Stability Verification}

A key robustness concern is whether \CRITIC weights are stable across the 14-year observation window. Substantial drift in weights would suggest either data quality issues or genuine structural breaks (regime changes) in governance dynamics, either of which would compromise the use of a single fixed weight vector.

A \textbf{split-half temporal stability test} is conducted:

1. \textbf{Partition data:} The 14-year panel is split at the temporal midpoint: Early period (2011--2017, $T_{\text{early}} = 7$ years) and Late period (2018--2024, $T_{\text{late}} = 7$ years).

2. \textbf{Compute weights independently:} The complete two-level \CRITIC pipeline is applied to each subsample separately. This yields two sets of 29 global weights: $\mathbf{w}_{\text{early}} = (w_1^{\text{early}}, \ldots, w_{29}^{\text{early}})^\top$ and $\mathbf{w}_{\text{late}}$.

3. \textbf{Compute cosine similarity:} The angle between the two weight vectors is measured using normalized cosine similarity:

\begin{equation}
\label{eq:cosine_sim}
\text{cos}(\theta) = \frac{\mathbf{w}_{\text{early}} \cdot \mathbf{w}_{\text{late}}}{\|\mathbf{w}_{\text{early}}\| \cdot \|\mathbf{w}_{\text{late}}\|}
\end{equation}

The cosine similarity lies in $[-1, 1]$, with $\text{cos}(\theta) = 1$ indicating perfect alignment, $\text{cos}(\theta) = 0$ indicating orthogonality, and $\text{cos}(\theta) = -1$ indicating opposite directions.

4. \textbf{Decision rule:} Weights are declared stable if:

\begin{equation}
\text{cos}(\theta) \geq \theta_{\text{stability}} \quad \text{(default: } \theta_{\text{stability}} = 0.95 \text{)}
\end{equation}

A threshold of 0.95 is sufficiently stringent to detect substantial weight drift while accommodating minor numerical fluctuations due to sample variation.

\textbf{Empirical result:} For the \PAPI dataset, $\text{cos}(\theta) = 0.9915$ is obtained, well above the 0.95 threshold, indicating very high temporal stability. This provides strong empirical support for treating the weight vector as a time-invariant global parameter suitable for ranking all years 2011--2024 under a common weighting scheme.

\subsubsection{Interpretation of Two-Level CRITIC Results}

The two-level framework yields interpretable results: High-variance, uncorrelated sub-criteria rank highly within their criterion group. Criteria with high variance and low inter-criterion correlation rank highly globally. In the empirical results section (Section~\ref{subsec:results_weights}), we report the 29 global weights, 8 criterion weights, and underlying variance\slash correlation diagnostics.

\subsection{MCDM Ranking Methods}
\label{subsec:mcdm}

\subsubsection{TOPSIS}
\label{subsubsec:topsis}
% ----------------------------------------------------------------
% PLACEHOLDER — Hwang & Yoon (1981). Closeness coefficient formula.
% ----------------------------------------------------------------
\textit{[Content placeholder.]}

\subsubsection{VIKOR}
\label{subsubsec:vikor}
% ----------------------------------------------------------------
% PLACEHOLDER — Opricovic & Tzeng (2004). S, R, Q indices. v = 0.5.
% ----------------------------------------------------------------
\textit{[Content placeholder.]}

\subsubsection{PROMETHEE II}
\label{subsubsec:promethee}
% ----------------------------------------------------------------
% PLACEHOLDER — Brans & Vincke (1985). V-shape preference function.
% Net flow phi_net.
% ----------------------------------------------------------------
\textit{[Content placeholder.]}

\subsubsection{COPRAS}
\label{subsubsec:copras}
% ----------------------------------------------------------------
% PLACEHOLDER — Zavadskas et al. (1994). S+, S-, Q formula.
% ----------------------------------------------------------------
\textit{[Content placeholder.]}

\subsubsection{EDAS}
\label{subsubsec:edas}
% ----------------------------------------------------------------
% PLACEHOLDER — Keshavarz Ghorabaee et al. (2015). PDA, NDA, AS formula.
% ----------------------------------------------------------------
\textit{[Content placeholder.]}

\subsubsection{Simple Additive Weighting (SAW)}
\label{subsubsec:saw}
% ----------------------------------------------------------------
% PLACEHOLDER — Fishburn (1967). Score_i = sum_j w_j * r_{ij}.
% ----------------------------------------------------------------
\textit{[Content placeholder.]}


\subsection{Temporal Feature Engineering}
\label{subsec:features}
% ----------------------------------------------------------------
% PLACEHOLDER — 12-block feature pipeline (~295 features per province):
%   Block 1: Current criterion values (C01–C08)
%   Block 2: Lag features (t-1, t-2, t-3) + _was_missing flags
%   Block 3: Rolling statistics (mean, std, min, max; windows 2, 3, 5)
%   Block 4: Momentum and acceleration
%   Block 5: Stationarity transformations (entity-demeaned, Δ², demeaned momentum)
%   Block 6: Polynomial trend slopes (≥3 valid pts)
%   Block 7: EWMA levels (spans 2, 3, 5)
%   Block 8: Expanding window mean
%   Block 9: Inter-criterion diversity (std, range)
%   Block 10: Rolling skewness (5-year window)
%   Block 11: Panel-relative features (percentile, z-score, rank-change)
%   Block 12: Geographic cluster dummies (5 Vietnam regions)
% Two-track preprocessing:
%   PLS track (BayesianRidge, KRR, SVR): PLSRegression(n_components=20)
%     with mutual-information pre-filter
%   Threshold track (CatBoost, QRF): variance threshold only
% ----------------------------------------------------------------
\textit{[Content placeholder.]}

\subsection{Base Forecasting Models}
\label{subsec:base_models}
% ----------------------------------------------------------------
% PLACEHOLDER — 5 diverse base learners:
%   1. CatBoost (joint multi-output MultiRMSE, oblivious trees)
%   2. Bayesian Ridge (PLS-compressed features, posterior uncertainty)
%   3. Quantile Random Forest (full predictive distributions)
%   4. Kernel Ridge Regression (RBF kernel, L2 regularisation)
%   5. Support Vector Regression (epsilon-insensitive tube, RBF)
% Target transform: logit for SAW [0,1] targets; Yeo-Johnson otherwise.
% ----------------------------------------------------------------
\textit{[Content placeholder.]}

\subsection{Super Learner Meta-Ensemble}
\label{subsec:superlearner}
% ----------------------------------------------------------------
% PLACEHOLDER — van der Laan et al. (2007) Super Learner.
% Panel-aware walk-forward CV (min_train_years = 8; 5 folds).
% Out-of-fold (OOF) predictions for each base model.
% Meta-learner: NNLS Ridge per output -> _meta_weights_per_output_ (n_outputs x n_models).
% Oracle inequality: Super Learner >= best convex combination asymptotically.
% OOF cache reused for conformal calibration (no re-fitting).
% Alternative meta-learners: Dirichlet stacking, Bayesian stacking, ElasticNet.
% ----------------------------------------------------------------
\textit{[Content placeholder.]}

\subsection{Conformal Prediction Intervals}
\label{subsec:conformal}
% ----------------------------------------------------------------
% PLACEHOLDER — Vovk et al. (2005), Angelopoulos & Bates (2022).
% Distribution-free 95% coverage guarantee (alpha = 0.05).
% CV+ conformal: calibration residuals from OOF predictions.
% Per-column OOF masks (_oof_valid_mask_per_col_) for multi-output.
% ConformalPredictor with calibrate_residuals(); warning at n_cal < 30.
% Alternatives: split conformal, Adaptive Conformal Inference (ACI).
% ----------------------------------------------------------------
\textit{[Content placeholder.]}

\subsection{Evaluation Metrics and Validation}
\label{subsec:evaluation}
% ----------------------------------------------------------------
% PLACEHOLDER — Forecast evaluation metrics: R², RMSE, MAE, MedAE, MAPE,
% Max Error, Bias; Winkler score for prediction interval quality.
% Residual diagnostics: Durbin-Watson (autocorrelation), Breusch-Pagan
%   (heteroscedasticity), Shapiro-Wilk (normality).
% Holdout evaluation: 2024 withheld, 252 OOF samples (genuinely out-of-sample).
% Sensitivity analysis: 1000 Monte Carlo simulations; weight perturbation ±10%;
%   Kendall rank correlation threshold = 0.85.
% ----------------------------------------------------------------
\textit{[Content placeholder.]}

% =============================================================
\section{Results}
\label{sec:results}
% =============================================================

\subsection{CRITIC Criterion Weights and Temporal Stability}
\label{subsec:results_weights}
% ----------------------------------------------------------------
% PLACEHOLDER — Report CRITIC criterion weights:
%   C05 (Public Admin Procedures): 19.62%
%   C03 (Vertical Accountability):  16.17%
%   C06 (Public Service Delivery):  15.61%
%   C04 (Control of Corruption):   12.96%
%   C01 (Participation):           12.72%
%   C02 (Transparency):            12.09%
%   C07 (Environmental Governance):  5.40%
%   C08 (E-Governance):              5.43%
% Top SC by global weight: SC32 (0.0730), SC54 (0.0539), SC53 (0.0514)
% Temporal stability cosine = 0.9915 (> 0.95 threshold).
% Table 2: Full 29 SC global weights (sorted by rank).
% Figure: weight radar chart.
% ----------------------------------------------------------------
\textit{[Content placeholder.]}

\subsection{Historical Provincial Rankings (2011--2024)}
\label{subsec:results_historical}
% ----------------------------------------------------------------
% PLACEHOLDER — Present Hierarchical MCDM results for all years.
% Kendall's W ~ 0.88 (strong inter-method agreement).
% Per-criterion MCDM score comparisons (TOPSIS, VIKOR, PROMETHEE II,
%   COPRAS, EDAS, SAW).
% Figure: parallel coordinate plot of per-criterion scores.
% Table 3: Top 10 provinces by SAW composite score, latest year (2024).
% Trend analysis: provinces improving/declining over 14 years.
% ----------------------------------------------------------------
\textit{[Content placeholder.]}

\subsection{Inter-Method Concordance and Robustness}
\label{subsec:results_concordance}
% ----------------------------------------------------------------
% PLACEHOLDER — Kendall's W concordance across 6 MCDM methods.
% Sensitivity analysis: 100+ Monte Carlo simulations with ±15%
%   weight perturbations. Robustness index and confidence intervals.
% Spearman rank correlations across methods.
% ----------------------------------------------------------------
\textit{[Content placeholder.]}

\subsection{Forecasting Model Performance}
\label{subsec:results_forecast_performance}
% ----------------------------------------------------------------
% PLACEHOLDER — 5-fold walk-forward CV R² per model:
%   BayesianRidge: 0.1304 ± 0.0514 (best individual)
%   KernelRidge:   0.1213 ± 0.0365 (most stable)
%   CatBoost:      0.1008 ± 0.0621
%   SVR:           0.0757 ± 0.0313
%   QuantileRF:   -0.0873 ± 0.1104
% Super Learner meta-weights:
%   BayesianRidge 49.0%, KernelRidge 19.1%, CatBoost 11.7%,
%   SVR 8.6%, QuantileRF 8.6%
% Holdout (2024 withheld, 252 OOF samples):
%   R² = 0.2171, RMSE = 2.3936, MAE = 1.1700
% Residual diagnostics: Durbin-Watson, Breusch-Pagan, Shapiro-Wilk.
% Figure: CV boxplots, feature importance heatmap, model weight chart.
% ----------------------------------------------------------------
\textit{[Content placeholder.]}

\subsection{2025 Provincial Governance Forecasts and Prediction Intervals}
\label{subsec:results_forecast_2025}
% ----------------------------------------------------------------
% PLACEHOLDER — 63-province 2025 forecasts for 8 criteria (SAW-normalised [0,1]).
% 95% conformal prediction intervals for each province-criterion cell.
% Interval width analysis: tightest vs. widest intervals by criterion.
% Table 4: Top 10 provinces in 2025 forecast ranking.
% Figure: prediction interval chart for selected provinces.
% Figure: conformal coverage calibration curve.
% ----------------------------------------------------------------
\textit{[Content placeholder.]}

\subsection{Integrated MCDM Ranking of 2025 Forecasts}
\label{subsec:results_2025_ranking}
% ----------------------------------------------------------------
% PLACEHOLDER — Apply CRITIC weights and 6 MCDM methods to 2025 forecasts.
% Produce 2025 forecast-based provincial ranking.
% Compare with 2024 historical ranking: rank shifts.
% Table 5: Full 63-province 2025 MCDM ranking and rank-change vs. 2024.
% ----------------------------------------------------------------
\textit{[Content placeholder.]}

% =============================================================
\section{Discussion}
\label{sec:discussion}
% =============================================================

\subsection{Interpretation of CRITIC Weights}
\label{subsec:disc_weights}
% ----------------------------------------------------------------
% PLACEHOLDER — Explain why C05 (Public Admin Procedures, 19.6%) receives
% the highest weight: high inter-provincial variance and low correlation
% with other criteria. Contrast with C07/C08 (low weights due to structural
% gap in early years reducing effective variance). Policy implications:
% administrative reform has highest leverage for differentiation.
% ----------------------------------------------------------------
\textit{[Content placeholder.]}

\subsection{MCDM Method Convergence and Governance Interpretation}
\label{subsec:disc_mcdm}
% ----------------------------------------------------------------
% PLACEHOLDER — W ≈ 0.88 confirms rankings are stable across methodological
% choices. Discuss cases where methods diverge (e.g., VIKOR vs TOPSIS in
% mid-ranking provinces) and how ER handles residual disagreement.
% Interpret top/bottom provinces. Connection to Vietnam Gov't policy targets.
% ----------------------------------------------------------------
\textit{[Content placeholder.]}

\subsection{Forecasting Performance in Small-Panel Contexts}
\label{subsec:disc_forecast}
% ----------------------------------------------------------------
% PLACEHOLDER — Individual model R² (0.07–0.13) is modest but expected for
% highly persistent panel data with structural breaks (COVID-19, sub-criteria
% introduction). The Super Learner's NNLS combination allocates near-zero
% weight to QRF when uncertainty is high (tree methods can overfit small T), while
% Bayesian Ridge (49%) and KernelRidge (19%) reflect suitability of
% regularised linear methods for small-T panels. Compare with benchmark
% (random walk, AR(1)). Holdout R² = 0.22 represents genuine out-of-sample
% predictive power.
% ----------------------------------------------------------------
\textit{[Content placeholder.]}

\subsection{Conformal Prediction Interval Quality}
\label{subsec:disc_conformal}
% ----------------------------------------------------------------
% PLACEHOLDER — Interval width variation across criteria: wider for criteria
% with structural breaks (C07, C08 missing early years), narrower for stable
% criteria (C01, C02). Empirical coverage vs. nominal 95%. Comparison with
% naive ±1.96*residual_std approach. Winkler score analysis.
% ----------------------------------------------------------------
\textit{[Content placeholder.]}

\subsection{Policy Implications for Vietnam's Governance Reform}
\label{subsec:disc_policy}
% ----------------------------------------------------------------
% PLACEHOLDER — Three actionable insights:
% (1) Provinces consistently underperforming on SC32 (govt responsiveness)
%     should be prioritised for judicial and public complaint reform.
% (2) The large 2025 forecast uncertainty bands for C07/C08 reflect
%     volatile growth in environmental enforcement and digital government;
%     policy investment in these areas should be accompanied by mid-year
%     monitoring.
% (3) The CRITIC-identified importance of C05/C06 suggests that national-level
%     administrative simplification programmes (e.g., single-window systems)
%     yield the highest marginal improvement in composite scores.
% ----------------------------------------------------------------
\textit{[Content placeholder.]}

\subsection{Comparison with Related Approaches}
\label{subsec:disc_comparison}
% ----------------------------------------------------------------
% PLACEHOLDER — Contrast with:
% (a) World Bank WGI (country-level, no sub-national differentiation)
% (b) PCI (Vietnam business climate index — focuses on investment climate,
%     not citizen governance)
% (c) Standard MCDM papers using entropy or AHP for weighting
% (d) ARIMA/VAR-based forecasting of governance indicators
% Advantages: objective weighting, multi-method consensus, ML forecasting,
%             distribution-free intervals, reproducible open-source pipeline.
% ----------------------------------------------------------------
\textit{[Content placeholder.]}

\subsection{Limitations}
\label{subsec:disc_limitations}
% ----------------------------------------------------------------
% PLACEHOLDER —
% (1) Survey-based data susceptible to political economy biases and mode effects.
% (2) Small T (14 years) limits statistical power of walk-forward CV and
%     motivates parsimonious models; forecasting gains are modest.
% (3) MCDM methods assume all-benefit (monotone) criteria and a static
%     preference structure; non-compensatory preferences not captured.
% (4) Conformal calibration pool (n_cal < 50 for some criteria-strata) may
%     yield conservative intervals under Mondrian stratification.
% (5) Causal interpretation of regression coefficients is not warranted;
%     the framework is predictive, not structural.
% ----------------------------------------------------------------
\textit{[Content placeholder.]}

% =============================================================
\section{Conclusion}
\label{sec:conclusion}
% =============================================================
% ----------------------------------------------------------------
% PLACEHOLDER — Summary of contributions:
% (1) First application of two-level CRITIC weighting to PAPI data;
%     confirmed temporal stability (cosine 0.9915).
% (2) Six-method MCDM consensus via ER; Kendall's W ≈ 0.88.
% (3) Three-tier ensemble (6 base models + Super Learner + conformal) for
%     governance forecasting; holdout R² = 0.22; 95% conformal intervals.
% (4) First forward-looking PAPI ranking for 2025.
%
% Limitations recap and future work:
% - Extend to district/commune level as granular data become available.
% - Incorporate text-based sentiment from citizen complaint records.
% - Apply Phase II enhancements: CQR, Mondrian CP, group LASSO meta-weights.
% - Investigate causal mechanisms (e.g., fiscal transfers → C03/C04 improvement).
% ----------------------------------------------------------------
\textit{[Content placeholder.]}

% =============================================================
%  REFERENCES
% =============================================================

\bibliography{references}

% INLINE BIBLIOGRAPHY — replace with BibTeX file or use \begin{thebibliography}
% The entries below are provided as a reference list for the BibTeX file.
% Create references.bib in the same directory with the entries listed below.

\begin{thebibliography}{99}

\bibitem[Acemoglu and Robinson(2005)]{Acemoglu2005}
Acemoglu, D., \& Robinson, J.~A. (2005).
\textit{Economic Origins of Dictatorship and Democracy}.
Cambridge University Press.

\bibitem[Angelopoulos and Bates(2022)]{Angelopoulos2022}
Angelopoulos, A.~N., \& Bates, S. (2022).
A gentle introduction to conformal prediction and
distribution-free uncertainty quantification.
\textit{arXiv preprint}, arXiv:2107.07511v7.

\bibitem[Brans and Vincke(1985)]{Brans1985}
Brans, J.~P., \& Vincke, P. (1985).
A preference ranking organization method (PROMETHEE).
\textit{Management Science}, 31(6), 647--656.

\bibitem[CECODES et al.(2024)]{PAPI2024}
CECODES, VFF-CRT, RTA, \& UNDP. (2024).
\textit{The 2024 Viet Nam Governance and Public Administration Performance Index
(PAPI): Measuring Citizens' Experiences}.
United Nations Development Programme, Ha Noi, Viet Nam.

\bibitem[Diakoulaki et al.(1995)]{Diakoulaki1995}
Diakoulaki, D., Mavrotas, G., \& Papayannakis, L. (1995).
Determining objective weights in multiple criteria problems:
The CRITIC method.
\textit{Computers \& Operations Research}, 22(7), 763--770.

\bibitem[Fishburn(1967)]{Fishburn1967}
Fishburn, P.~C. (1967).
Additive utilities with incomplete product sets:
Application to priorities and assignments.
\textit{Operations Research}, 15(3), 537--542.

\bibitem[Greco et al.(2019)]{Greco2019}
Greco, S., Ishizaka, A., Tasiou, M., \& Torrisi, G. (2019).
On the methodological framework of composite indices: A review
of the issues of weighting, aggregation, and robustness.
\textit{Social Indicators Research}, 141(1), 61--94.

\bibitem[Hwang and Yoon(1981)]{Hwang1981}
Hwang, C.~L., \& Yoon, K. (1981).
\textit{Multiple Attribute Decision Making: Methods and Applications}.
Springer-Verlag, Berlin.

\bibitem[Kaufmann et al.(2010)]{Kaufmann2010}
Kaufmann, D., Kraay, A., \& Mastruzzi, M. (2010).
The worldwide governance indicators: Methodology and analytical issues.
\textit{Hague Journal for the Rule of Law}, 3(2), 220--246.

\bibitem[Keshavarz~Ghorabaee et al.(2015)]{Keshavarz2015}
Keshavarz~Ghorabaee, M., Zavadskas, E.~K., Olfat, L., \& Turskis, Z. (2015).
Multi-criteria inventory classification using a new method of evaluation
based on distance from average solution (EDAS).
\textit{Informatica}, 26(3), 435--451.

\bibitem[Malesky(2014)]{Malesky2014}
Malesky, E.~J. (2014).
Rethinking the relationship between subnational government and economic
performance: Evidence from Vietnam.
\textit{Journal of East Asian Studies}, 14(3), 399--452.

\bibitem[Nguyen et al.(2020)]{Nguyen2020}
Nguyen, H.~T., Pham, T.~H., \& Vu, T.~A. (2020).
Decentralisation, governance quality, and economic growth in Vietnam:
A provincial-level panel study.
\textit{Journal of Development Economics}, 144, 102436.

\bibitem[Opricovic and Tzeng(2004)]{Opricovic2004}
Opricovic, S., \& Tzeng, G.~H. (2004).
Compromise solution by MCDM methods: A comparative analysis of
VIKOR and TOPSIS.
\textit{European Journal of Operational Research}, 156(2), 445--455.

\bibitem[Prokhorenkova et al.(2018)]{Prokhorenkova2018}
Prokhorenkova, L., Gusev, G., Vorobev, A., Dorogush, A.~V., \& Gulin, A. (2018).
CatBoost: Unbiased boosting with categorical features.
\textit{Advances in Neural Information Processing Systems}, 31.

\bibitem[van~der~Laan et al.(2007)]{vanderLaan2007}
van~der~Laan, M.~J., Polley, E.~C., \& Hubbard, A.~E. (2007).
Super learner.
\textit{Statistical Applications in Genetics and Molecular Biology}, 6(1),
Article~25.

\bibitem[Vovk et al.(2005)]{Vovk2005}
Vovk, V., Gammerman, A., \& Shafer, G. (2005).
\textit{Algorithmic Learning in a Random World}.
Springer, New York.

\bibitem[Yang and Xu(2002)]{Yang2002}
Yang, J.~B., \& Xu, D.~L. (2002).
On the evidential reasoning algorithm for multiple attribute decision analysis
under uncertainty.
\textit{IEEE Transactions on Systems, Man, and Cybernetics---Part~A:
Systems and Humans}, 32(3), 289--304.

\bibitem[Zavadskas et al.(1994)]{Zavadskas1994}
Zavadskas, E.~K., Kaklauskas, A., \& Sarka, V. (1994).
The new method of multicriteria complex proportional assessment of projects.
\textit{Technological and Economic Development of Economy}, 1(3), 131--139.

\bibitem[Doove et al.(2014)]{Doove2014}
Doove, L.~L., van Buuren, S., \& Dusseldorp, E. (2014).
Recursive partitioning for missing data imputation in the presence of
non-linear relations.
\textit{Computational Statistics \& Data Analysis}, 72, 92--104.

\bibitem[Little(1988)]{Little1988}
Little, R.~J.~A. (1988).
A test of missing completely at random for multivariate data with missing values.
\textit{Journal of the American Statistical Association}, 83(404), 1198--1202.

\bibitem[Rubin(1987)]{Rubin1987}
Rubin, D.~B. (1987).
\textit{Multiple Imputation for Nonresponse in Surveys}.
Wiley, New York.

\bibitem[Scikit-learn Developers(2023)]{scikit-learn}
Scikit-learn Developers (2023).
Scikit-learn: Machine learning in Python.
\textit{Journal of Machine Learning Research}, 12, 2825--2830.

\bibitem[Sterne et al.(2009)]{Sterne2009}
Sterne, J.~A.~C., White, I.~R., Carlin, J.~B., Spratt, M., Royston, P.,
Kenward, M.~G., Wood, A.~M., \& Carpenter, J.~R. (2009).
Multiple imputation for missing data in epidemiological and clinical research:
Potential and pitfalls.
\textit{BMJ}, 338, b2393.

\bibitem[van Buuren and Groothuis-Oudshoorn(2011)]{vanBuuren2011}
van Buuren, S., \& Groothuis-Oudshoorn, K. (2011).
\texttt{mice}: A package for creating multiple imputations by chained equations.
\textit{Journal of Statistical Software}, 45(3), 1--67.

\bibitem[van Buuren(2018)]{vanBuuren2018}
van Buuren, S. (2018).
\textit{Flexible Imputation of Missing Data} (2nd ed.).
Chapman and Hall/CRC, Boca Raton, FL.

\end{thebibliography}

\end{document}
